\documentclass[11pt,oneside]{article}

\usepackage[margin=1.15in]{geometry}
\usepackage{setspace}
\setstretch{1.15}

\usepackage{fontspec}
\usepackage{unicode-math}
\setmainfont{Libertinus Serif}
\setsansfont{Libertinus Sans}
\setmonofont{Libertinus Mono}
\setmathfont{Libertinus Math}

\usepackage{amsmath,amssymb}
\usepackage{microtype}
\usepackage{hyperref}

\title{Induced Homogeneity\\
{\normalsize Aspirational Pages and the Collapse of Meaning Under Engagement Optimization} }
\author{Flyxion}
\date{\today}

\begin{document}
\maketitle

\begin{abstract}
Contemporary social platforms exhibit a pervasive proliferation of minimally active Pages bearing aspirational identities yet lacking sustained practice. This phenomenon is often misattributed to coordinated content farms or malicious automation. This paper advances an alternative account: that such Pages emerge from tutorialized onboarding architectures that lower the barrier to identity declaration while failing to enforce capability acquisition. When platforms invite users to name, brand, or position themselves prior to demonstrating practice, aspiration becomes uncoupled from action, producing large populations of ontologically incomplete entities.

To formalize this failure mode, the paper distinguishes axiomatically between aspiration, identity, and practice, and proves that identity stabilization requires persistence of namespace binding, append-only provenance, and bounded semantic drift. Using the Relativistic Scalar--Vector--Entropy Plenum (RSVP) framework, the dynamics of stalled Pages are modeled as regions of elevated semantic potential without corresponding vector realization, leading to entropy accumulation. A decay law is derived for unsupported identity claims, and it is shown that downstream moderation cannot compensate for misordered onboarding.

Counterexamples drawn from successful instructional systems, including capability-scaffolded software tutorials and embodied demonstrations of practice, show that tutorials can enforce action before recognition and thereby conserve meaning. The paper concludes by translating these insights into platform design constraints, arguing that identity must emerge as a downstream effect of demonstrated practice rather than as an entry condition. Meaning, it is argued, cannot be optimized into existence through engagement incentives, but must be conserved through architectural coupling between action, history, and identity.
\end{abstract}


\section{Introduction}

The contemporary proliferation of minimally active Facebook Pages bearing aspirational titles is often mischaracterized as the product of coordinated content farms or malicious automation. While such actors exist, this explanation fails to account for the far larger and more revealing pattern: the convergence of thousands of independent individuals on nearly identical semantic forms. These Pages are not primarily produced by centralized organizations, but by ordinary users responding to a platform that has aggressively lowered the barrier to entry for Page creation while simultaneously withholding the structural conditions required for meaningful participation.

What appears, from a distance, to be an industrial content farm is in fact an emergent phenomenon produced by tutorialized aspiration. The platform teaches users how to enter the game, but not how to inhabit it.

\section{The Algorithmic Tutorial as Cultural Force}

Facebook’s Page-creation flows are no longer neutral tools. They function as implicit tutorials that frame participation as an immediately monetizable or reputation-generating activity. Users are encouraged to create Pages as a first step toward visibility, income, or recognition, often before they possess the social, experiential, or expressive resources necessary to sustain original content production.

The tutorial does not instruct users in craft, voice, or domain mastery. It instructs them in naming, categorization, and basic posting mechanics. As a result, users select Page names that feel maximally promising within the narrow symbolic vocabulary presented to them. Aspirational titles referencing professions, lifestyles, or achievement thresholds become default choices not because they are strategic, but because they are legible within the tutorial’s implied success narrative.

\section{Aspirational Naming as Pre-Identity}

The resulting Page names function as pre-identities. They are projections of a hoped-for future state rather than descriptions of an existing practice. The Page is created in advance of the life it is meant to represent. In many cases, the user intends to grow into the name, acquiring skills, experience, and audience after the fact.

This temporal inversion is crucial. Identity precedes substance. The Page exists before the practice that would justify it. When the anticipated growth fails to materialize, the Page remains as an inert artifact: not fraudulent, but unfinished.

\section{Skill Asymmetry and Expressive Friction}

Most participants in this system are not deceptive. They are inexperienced. Producing sustained original content requires competencies that are unevenly distributed and slow to acquire: narrative coherence, technical skill, aesthetic judgment, and an understanding of audience dynamics. The platform, however, treats these as optional. Its interfaces suggest that content production is trivial, repeatable, and immediately rewarded.

When users encounter the actual difficulty of creation, many stall. The Page remains nominally active but substantively empty. The user does not necessarily abandon the project out of bad faith, but out of mismatch between expectation and capacity.

\section{Emergent Homogeneity Without Coordination}

Because the tutorial, incentives, and symbolic vocabulary are shared, the resulting Pages resemble one another even though they are independently created. Homogeneity emerges without coordination. The platform produces a statistical aesthetic: similar names, similar categories, similar stalled trajectories.

From the outside, this resembles a content farm. From the inside, it is a mass of individuals making locally rational choices under identical constraints and affordances. The system does not require bots to generate semantic noise; it recruits human aspiration itself as the generator.

\section{Engagement Optimization and the Indifference to Completion}

The engagement-optimized architecture of the platform is indifferent to whether these aspirational projects mature. A Page that never posts meaningfully does not significantly harm platform metrics. Its mere existence expands the content surface area and sustains the narrative of opportunity. The system rewards entry, not completion.

This indifference is structurally significant. In systems that conserve meaning, incomplete entities decay or consolidate. Here, they persist indefinitely, accumulating as semantic residue.

\section{Identity Without Cost}

The underlying failure is the absence of cost attached to identity creation and abandonment. Pages can be created speculatively, left dormant, and replaced without consequence. Identity does not bind action to history. Names do not obligate performance. Time does not accumulate as responsibility.

Within your framework of identity as namespace, this represents a breakdown of attributional constraint. Without persistent binding between name, action, and consequence, aspiration decouples from accountability and meaning dissipates into noise.

\section{Entropy and the Tutorial Trap}

The tutorialized entry system introduces entropy faster than the platform can dissipate it. Each new Page adds semantic mass without corresponding structure. The system lacks mechanisms to enforce maturation, collapse unfinished attempts, or reclaim abandoned identity space.

This produces what might be called a tutorial trap: users are invited to play, but the rules of meaningful play are never encoded into the system. Optimization operates downstream of incoherence.

\section{Toward Constraint-First Remediation}

A coherence-first architecture would reverse this dynamic. Identity creation would be gated not by verification alone, but by demonstrated continuity of practice. Names would bind behavior across time, and aspirational projection would incur informational debt rather than being cost-free. Pages would be required to evolve through traceable transformations rather than appearing fully formed at the outset.

Such constraints would not suppress participation. They would align participation with reality. Users would still aspire, but aspiration would be coupled to process rather than tutorial completion.

\section{Conclusion}

The proliferation of hollow aspirational Pages on Facebook is not primarily a story of manipulation or centralized exploitation. It is a story of people responding sincerely to an environment that promises immediacy while withholding structure. The platform teaches users how to enter, but not how to become.

What emerges is not fraud, but entropy. Meaning fragments because identity is unbound, history is optional, and aspiration is allowed to precede competence without consequence. Until platforms treat identity, continuity, and semantic cost as foundational constraints, they will continue to generate landscapes that look industrially farmed even when they are populated by individuals acting in good faith.

Meaning cannot be crowdsourced into existence through tutorials. It must be scaffolded by architecture.

\section{Axiomatic Distinction Between Aspiration, Identity, and Practice}

To understand why tutorial-induced participation produces semantic residue rather than durable meaning, it is necessary to distinguish formally between aspiration, identity, and practice. These are often conflated in platform design, yet they occupy fundamentally different roles in the construction of meaning-bearing systems.

Aspiration refers to a projected future state toward which an agent or entity intends to move. It is prospective, informationally light, and unconstrained by evidence. Aspiration is not false by default, but it is unverified. Its function is motivational rather than descriptive.

Identity, in contrast, is a binding structure. It functions as a namespace that associates actions, representations, and consequences across time. An identity is meaningful only insofar as it accumulates history and constrains future transformation. Identity is not merely a label but a persistent reference frame that allows observers and systems to attribute continuity.

Practice is the realized process through which aspiration is tested and identity is substantiated. It consists of repeated, situated actions that generate evidence. Practice is costly in time, effort, and exposure to failure, and it is through this cost that meaning is stabilized.

The failure mode observed on engagement-optimized platforms arises when identity is instantiated at the level of aspiration rather than practice. The system allows a projected future state to be named as if it were already realized. The Page name functions as an identity claim without requiring any corresponding accumulation of practice. When practice fails to materialize, the identity does not collapse or transform; it simply persists in an inert state.

This inversion violates a necessary ordering constraint. In coherent systems, aspiration may precede practice, but identity must be downstream of demonstrated continuity. When identity is allowed to precede practice without cost, the system fills with unverified namespaces that carry no informational weight. Meaning becomes indistinguishable from intent, and attribution loses its grounding.

Within this axiomatic frame, the platform’s tutorial mechanisms can be seen as systematically encouraging premature identity instantiation. They reward naming before doing, projection before accumulation, and visibility before substantiation. The resulting Pages are not deceptive but ontologically incomplete. They are identities without proofs.

\section{RSVP Field Interpretation of Tutorial-Induced Entropy}

The dynamics described above can be rendered with greater precision by mapping them onto the Relativistic Scalar--Vector--Entropy Plenum framework. In this interpretation, meaning-bearing activity on a platform is modeled as a coupled field system in which scalar potential, directed flow, and entropy interact over time.

Aspiration corresponds to elevated scalar potential. It represents stored semantic energy: the promise of future structure. The platform’s tutorials artificially inflate this potential by encouraging users to instantiate identities that signal high future payoff without requiring commensurate field gradients.

Practice corresponds to vector flow. It is the directed movement of semantic material through time via repeated action. Sustained content creation, audience interaction, and refinement constitute coherent flow. When practice is absent or stalls, vector magnitude approaches zero even if scalar potential remains high.

Entropy accumulates when scalar potential is not discharged through flow. In the observed phenomenon, Pages with aspirational names possess elevated semantic potential but negligible vector realization. The system permits this configuration to persist indefinitely, resulting in localized entropy pockets: regions of stored promise that never resolve into structure.

In a healthy semantic field, such configurations are unstable. Excess potential without flow induces corrective dynamics. Either flow emerges, potential dissipates, or the structure collapses and releases its stored energy back into the surrounding field. Engagement-optimized platforms suppress these corrective mechanisms. They allow stalled high-potential nodes to persist without decay.

This suppression produces a characteristic field pathology. The semantic landscape becomes littered with shallow potential wells that attract attention weakly and inconsistently, disrupting coherent trajectories. Legitimate flows must navigate around these inert regions, increasing noise and reducing signal-to-noise ratios across the system.

From the RSVP perspective, the tutorial system acts as a scalar inflation mechanism without corresponding vector constraints. It creates promise without process. Entropy is not merely a side effect but the inevitable result of uncoupled field components.

A coherence-first architecture would reintroduce coupling. Scalar potential would be permitted to rise only in proportion to demonstrated vector flow. Identity claims would emerge gradually as integrated field structures rather than instant labels. Entropy would be conserved through enforced decay, consolidation, or reintegration of stalled regions.

Seen this way, the proliferation of hollow Pages is not an anomaly but a predictable phase state of a system that has broken the coupling between aspiration, action, and consequence. Until that coupling is restored, the field will continue to fill with ungrounded potential, and meaning will remain diffuse.

\section{A Proposition on Necessary Conditions for Identity Stabilization}

The preceding analysis treats identity not as a cosmetic label but as an infrastructural binding that permits attribution, accumulation, and constraint across time. In that sense, identity stabilization is not a psychological phenomenon but a systems property: it is the condition under which an entity remains legible as ``the same'' actor under successive interactions, while retaining sufficient semantic integrity that past behavior meaningfully constrains future expectation.

Let an entity $e$ be represented on a platform by a name or handle $n$ together with a history of observable actions $H_e(t)$ over time. Let $M_e(t)$ denote the semantic state attributed to $e$ by the system and its observers at time $t$, and let $I_e(t)$ denote the identity binding relation that maps actions into a persistent namespace. The platform induces an evolution operator $\mathcal{T}$ such that
\[
M_e(t+\Delta t) = \mathcal{T}\big(M_e(t), a_e(t), \Delta t; \Pi\big),
\]
where $a_e(t)$ is the action (or null action) emitted by $e$ at time $t$ and $\Pi$ denotes platform governance and recommendation parameters.

We require three structural ingredients for identity stabilization: persistence, traceability, and bounded semantic drift. Persistence means that the mapping from $e$ to its namespace does not reset cheaply. Traceability means that the semantic state is conditioned by an append-only transformation history rather than overwritten by redefinition. Bounded drift means that successive transformations cannot arbitrarily change the functional role of the entity without incurring cost or exhausting a budget.

These notions can be expressed through three constraints. First, there exists a binding operator $\mathcal{B}$ such that the namespace $I_e(t)$ is a function of the full history rather than a freely editable label, with stability expressed by a Lipschitz-like condition
\[
d_I\!\left(I_e(t+\Delta t), I_e(t)\right) \leq \kappa\, \|a_e(t)\| \Delta t,
\]
for some $\kappa>0$ and an identity metric $d_I$ that measures namespace mutation. This constraint captures the idea that identity can evolve, but only in proportion to realized action.

Second, there exists an append-only provenance map $\mathcal{P}$ such that
\[
H_e(t+\Delta t) = H_e(t) \oplus \mathcal{P}\!\left(a_e(t)\right),
\]
where $\oplus$ denotes irreversible accumulation. Without this property, the system cannot enforce accountability because history can be rewritten.

Third, there exists a semantic budget functional $B_e(t)$ with depletion under transformation, such that semantic drift satisfies
\[
d_M\!\left(M_e(t+\Delta t), M_e(t)\right) \leq B_e(t) - B_e(t+\Delta t),
\]
where $d_M$ is a metric on semantic states. This captures the claim that meaning is conserved only if changes spend a scarce resource.

The relevant claim can then be stated as a proposition.

\textbf{Proposition (Necessary Conditions for Identity Stabilization).} Consider a platform in which entities are represented by namespaces with attributed semantic states $M_e(t)$ and binding relations $I_e(t)$. If any of the following conditions fails, identity stabilization cannot hold at scale: (i) the namespace binding operator is persistent in the sense that namespace change is strictly coupled to realized action; (ii) the provenance map is append-only and publicly auditable at the level required for attribution; (iii) semantic drift is bounded by a scarce budget that cannot be replenished faster than it is spent through verifiable practice. Conversely, if all three conditions hold and the platform’s amplification operator is conditioned on accumulated provenance rather than instantaneous engagement, then there exists a regime in which identities become asymptotically stable, in the sense that observer-attributed semantic variance decreases with time under continued practice.

\textit{Sketch of justification.} If persistence fails, identities can be regenerated or rebranded without cost, producing attributional aliasing and eliminating the informational value of history. If provenance fails, history cannot constrain expectation because it is not reliably conserved. If bounded drift fails, semantic roles can change without cost, allowing the same namespace to become arbitrarily many functional entities over time, which is equivalent to identity collapse. When all three constraints hold, repeated practice increases the density of binding evidence, and the effective posterior uncertainty over an entity’s semantic role decreases. In that regime, identities stabilize as an emergent consequence of conserved history plus constrained transformation.

\section{A Decay Law for Stalled Pages in RSVP Terms}

The ``stalled Page'' phenomenon is naturally expressed in RSVP as a mismatch between scalar potential and vector realization. Let $\Phi_e(t)$ denote the scalar semantic potential associated with entity $e$, interpreted as accumulated promise, projected identity, or attention-attracting plausibility. Let $\mathbf{v}_e(t)$ denote the vector of realized practice, interpreted as directed, sustained activity producing verifiable transformations. Let $S_e(t)$ denote the local semantic entropy, interpreted as the degree of incoherence, ambiguity, and unresolvable attribution generated by the entity’s presence.

A stalled Page is characterized by persistent $\Phi_e(t)$ with negligible $\|\mathbf{v}_e(t)\|$. In a coherence-preserving system, such a configuration should not be metastable. It should either convert potential into flow (practice emerges), or potential should dissipate and the entity should decay, collapsing back into the background namespace.

We can model this with a dissipation equation in which potential decays when it is not supported by flow. A minimal decay law is
\[
\frac{d\Phi_e}{dt} = -\alpha\,\Phi_e + \beta\,\|\mathbf{v}_e\| - \gamma\,\Phi_e\,\sigma\!\left(\frac{\Phi_e}{\Phi_0}\right)\left(1-\sigma\!\left(\frac{\|\mathbf{v}_e\|}{v_0}\right)\right),
\]
where $\alpha>0$ is baseline dissipation, $\beta\ge 0$ is reinforcement by realized practice, and the final term represents nonlinear decay when potential is high but flow is low. The function $\sigma(x)=1/(1+e^{-x})$ is a smooth thresholding function, and $\Phi_0, v_0$ define scale parameters.

The interpretation is straightforward. Potential decays slowly by default, but is replenished when practice is present. When potential is large and flow remains below a threshold, the system applies accelerated decay. This captures the intuition that ambitious identity claims without substantiation should not persist indefinitely.

Entropy dynamics follow from uncoupled potential. A minimal coupling is
\[
\frac{dS_e}{dt} = \eta\,\Phi_e\left(1-\sigma\!\left(\frac{\|\mathbf{v}_e\|}{v_0}\right)\right) - \lambda\,\|\mathbf{v}_e\| + D_S\,\Delta S_e,
\]
where $\eta>0$ converts unsupported potential into entropy, $\lambda>0$ represents entropy reduction through practice that resolves ambiguity, and $D_S$ is a diffusion coefficient capturing how incoherence spreads into the surrounding semantic medium.

A system that currently tolerates massive numbers of inert Pages behaves as though $\alpha$ is near zero, the nonlinear decay term is absent or suppressed, and the amplification system injects external potential that is not conditional on flow. In RSVP terms, it is a platform that inflates $\Phi$ without enforcing a corresponding $\mathbf{v}$ constraint, thereby allowing $S$ to accumulate as a stable byproduct.

The practical consequence is that the semantic medium becomes cluttered with high-$\Phi$, low-$\mathbf{v}$ nodes that do not resolve into structure. The decay law above is therefore not merely descriptive; it is a candidate governance mechanism: an architectural reintegration process that converts unsubstantiated potential into decay rather than indefinite persistence.

\section{Translation into Platform Design Constraints}

The preceding proposition and decay law translate into concrete platform constraints when identity is treated as infrastructure and meaning is treated as conserved rather than extracted. The first constraint is that identity claims must be coupled to verifiable practice. This does not imply elitism or exclusion; it implies that the platform should not treat the mere creation of a Page name as equivalent to the existence of a meaning-bearing actor. The interface may allow early-stage experimentation, but the system must encode an ontological distinction between provisional aspiration and stabilized identity.

The second constraint is that histories must be append-only and legible to governance. A Page’s semantic state must be derived from transformation history rather than overwritten by cosmetic edits. Renaming, category shifts, and functional repurposing should be modeled as typed transformations with explicit provenance. In the absence of this structure, the platform cannot bind expectations to entities, and users cannot rationally allocate trust.

The third constraint is that semantic drift must be budgeted. If a Page wishes to change categories, branding, or functional role, it should spend a scarce resource that is replenished only by sustained practice. The budget is not a moral punishment but a thermodynamic accounting device: it prevents cheap redefinition from erasing accumulated meaning.

The fourth constraint is the introduction of controlled decay for stalled entities. Provisional Pages that exhibit high aspirational potential but negligible realized flow should not persist indefinitely as inert semantic mass. Instead, they should transition into a dormant state that is explicitly marked as such, with reduced algorithmic visibility and a gradual dissolution of representational privilege unless practice resumes. This is the design analog of the RSVP decay law: unsupported potential must dissipate.

The fifth constraint is that amplification must be conditioned on provenance density rather than instantaneous engagement. If recommendation rewards early spikes independent of continuity, it will incentivize premature identity instantiation and will fill the system with ungrounded claims. If recommendation instead treats continuity and typed transformation histories as primary signals, then stable identity becomes the easiest path to visibility, and tutorialized aspiration becomes less likely to crystallize into permanent debris.

Under these constraints, the platform remains open to novices while no longer mistaking entry for legitimacy. Participation is not suppressed; it is given a coherent developmental pathway. Aspiration is permitted, but it is not allowed to masquerade indefinitely as completed identity. In RSVP terms, the system enforces coupling between $\Phi$ and $\mathbf{v}$ and prevents $S$ from becoming a stable equilibrium state. Meaning, once again, becomes a conserved quantity that emerges from practice rather than from naming.

\section{Counterexamples: Tutorials That Produce Stabilized Competence}

The preceding critique should not be read as an indictment of tutorials as such. The failure mode described is not intrinsic to instructional scaffolding, but to tutorials that decouple identity projection from embodied practice. There exist well-established counterexamples in which tutorial systems succeed precisely because they refuse to grant symbolic identity or status in advance of demonstrated competence, and instead force meaning to emerge through action.

A canonical example is the built-in tutorial of the \texttt{vim} text editor. The tutorial does not ask the learner to declare themselves a programmer, an editor power-user, or a member of any aspirational category. It does not promise visibility, monetization, or recognition. Instead, it places the learner immediately inside the operational environment and constrains progress through concrete tasks. Each concept is introduced only insofar as it is exercised. Advancement is inseparable from practice, and the learner exits the tutorial not with a badge or a title, but with a usable skill that persists outside the tutorial context. Identity emerges downstream of capability rather than preceding it.

A related but more explicit instantiation of this principle is found in \textit{1}, a two-dimensional platformer designed to teach the same editing paradigm through embodied interaction. Here the tutorial is not merely instructional but structural. Movement, navigation, and success are impossible without executing the relevant commands. The system does not simulate learning; it enforces it. The player cannot meaningfully inhabit the game world without acquiring operational fluency, and the completion of the tutorial leaves the participant transformed in a way that is externally verifiable. The tutorial collapses into practice, and the game dissolves into a transferable competence.

An even earlier example can be found in the tutorial environment of Microsoft Windows 3.1. The system assumed no prior familiarity with graphical interfaces and therefore did not abstract interaction into symbolic explanation. Instead, it guided users through direct manipulation of windows, icons, and controls, asking them to perform actions rather than merely observe them. The tutorial embedded worked examples within the live environment and then required the user to reproduce those actions independently. Crucially, no persistent identity was created at the tutorial stage. There was no notion of becoming a “Windows user” in name before demonstrating basic operational mastery. Competence was established first; identity followed implicitly through use.

What unites these counterexamples is not polish or nostalgia, but structural discipline. In each case, aspiration is deliberately minimized. The tutorial does not invite the learner to imagine a future self defined by status or recognition. It invites them to act. Scalar potential is kept low and local. Vector flow is enforced from the outset. As a result, entropy does not accumulate. There are no large populations of stalled participants who have declared an identity without acquiring capacity, because the system does not permit identity to be declared at all.

In RSVP terms, these tutorials maintain tight coupling between scalar potential and vector realization. The system injects only as much potential as can be immediately discharged through action. Any failure to act results not in indefinite stasis, but in immediate friction. There is no stable equilibrium in which promise can persist without flow. Entropy is therefore continuously dissipated through practice rather than allowed to accumulate as semantic residue.

These examples clarify that the failure of contemporary platform tutorials is not inevitable. It is the result of a specific architectural choice: to treat naming, visibility, and monetization as entry conditions rather than as emergent properties of sustained practice. Where tutorials are designed to withhold identity until competence is demonstrated, they function as true onboarding mechanisms. Where tutorials instead invite users to claim identity in advance, they function as aspiration amplifiers and generate exactly the stalled, hollow entities observed at scale.

The lesson is not that platforms should abandon tutorials, but that tutorials must be ontologically honest. They must refuse to simulate arrival. They must force becoming.

\section{Tutorials as Capability Scaffolds versus Aspiration Inflators}

The contrast between successful instructional systems and contemporary platform onboarding can be formalized as a distinction between tutorials that scaffold capability and tutorials that inflate aspiration. Both are superficially similar in that they present themselves as entry points, yet they differ fundamentally in the ordering they impose between action, identity, and reward.

In a capability-scaffold tutorial, the system withholds symbolic recognition until operational competence is demonstrated. Progress is gated by action, not declaration. The learner is not invited to imagine a future self defined by status, but is instead constrained to operate within a narrow task environment where each step requires execution rather than affirmation. Identity emerges implicitly as a residue of accumulated practice. The tutorial collapses into the activity it teaches, and the boundary between learning and use becomes indistinct.

In contrast, aspiration-inflating tutorials foreground symbolic identity at the moment of entry. They invite users to name themselves, brand themselves, or position themselves within an imagined future of visibility or monetization before any sustained practice has occurred. The tutorial signals arrival rather than initiation. Progress is measured by completion of setup rituals rather than by demonstrated capability, and the system treats declaration as a sufficient proxy for competence. Practice, if it occurs, is deferred and optional.

This distinction explains why aspiration-inflating systems reliably produce large populations of stalled entities. By granting identity upstream of action, they permit scalar semantic potential to accumulate without enforcing vector realization. Users are encouraged to inhabit a symbolic role whose maintenance requires skills they have not yet acquired, and when the difficulty of practice becomes apparent, the system offers no corrective pressure. The identity remains instantiated even as action ceases, producing inert semantic mass.

The capability-scaffold model, by contrast, does not permit such stalling because it offers no stable configuration in which identity can persist without action. Failure to act results not in dormant existence but in immediate friction. The system does not shame or exclude; it simply refuses to stabilize meaning prematurely.

This distinction can be formalized as a corollary to the earlier proposition on identity stabilization.

\textbf{Corollary (Tutorial Ordering Constraint).} In any platform where tutorial mechanisms grant durable identity or visibility prior to the accumulation of verifiable practice, identity stabilization cannot occur at scale, even if persistence, provenance, and bounded drift are nominally enforced elsewhere in the system. Conversely, in platforms where tutorial mechanisms enforce action before identity recognition, identity stabilization emerges naturally as a downstream effect of accumulated practice, without requiring additional moderation or policing.

\textit{Justification.} Tutorials define the initial conditions of participation. If the initial condition allows identity instantiation without practice, then the binding operator is violated at the point of entry, and no downstream constraint can fully repair the resulting attributional entropy. The system begins with high scalar potential and zero vector flow, a configuration that admits stable stalled equilibria unless explicitly prohibited. If, however, the tutorial enforces practice as a prerequisite for recognition, then the binding operator is satisfied from the outset. Identity emerges only after sufficient evidence has accumulated, and the posterior uncertainty associated with the entity decreases monotonically with continued action.

In RSVP terms, the tutorial phase determines whether the system injects uncoupled scalar potential or immediately couples potential to flow. Aspiration-inflating tutorials act as scalar pumps, raising $\Phi$ without regard to $\mathbf{v}$ and thereby seeding entropy. Capability-scaffold tutorials act as coupling operators, ensuring that any increase in $\Phi$ is locally discharged through enforced flow. The difference is not pedagogical but ontological. One produces promise without structure; the other produces structure without promise, allowing promise to emerge later as a derived property.

This analysis clarifies why certain instructional systems scale gracefully while engagement-driven platforms do not. The issue is not that users lack motivation or talent, but that the architecture teaches them to declare before they can do. When tutorials simulate arrival rather than enforce becoming, identity ceases to function as infrastructure and degenerates into a decorative label. Meaning, once again, is lost not through malice or incompetence, but through misordered design.

\section{Empirical Signatures of Aspiration Inflation and Capability Scaffolding}

The distinction between aspiration-inflating tutorials and capability-scaffolding tutorials is not merely theoretical; it produces observable, system-level signatures that can be empirically distinguished. These signatures arise from differences in how semantic potential, realized practice, and entropy evolve over time within the platform.

In aspiration-inflating systems, one observes a sharp initial surge in entity creation accompanied by a rapid decay in sustained activity. The distribution of entities exhibits a heavy tail of near-zero practice histories coupled with persistent identity markers such as names, categories, or branding elements. Engagement, where present, is episodic and weakly correlated with continuity. Semantic states fluctuate without convergence, and rebranding or repurposing events occur without proportional accumulation of prior practice. Over time, the system’s content surface becomes increasingly cluttered with inert or weakly animated entities whose existence cannot be predicted from their activity traces.

In capability-scaffolding systems, by contrast, entity creation rates are lower but more strongly correlated with subsequent action density. Early abandonment is visible, but it manifests as immediate friction rather than long-lived dormancy. The distribution of practice histories shows fewer zero-length trajectories and a tighter coupling between identity persistence and accumulated action. Semantic roles converge rather than drift, and entities that persist exhibit decreasing variance in observer-attributed meaning. The system’s semantic medium remains sparse enough that discovery and attribution retain signal integrity.

These signatures are measurable without access to internal platform data. Publicly observable indicators such as posting cadence, rename frequency, category volatility, and lifespan-to-output ratios suffice to distinguish whether a platform’s onboarding architecture privileges aspiration or enforces capability. The presence of large populations of permanently stalled identities is therefore not a sociological mystery but an architectural fingerprint.

\section{A Design Theorem on the Limits of Moderation}

The preceding analysis permits a stronger formal claim regarding platform governance. Specifically, it allows one to state a theorem about the insufficiency of downstream moderation in systems whose onboarding order is misaligned.

\textbf{Theorem (Moderation Insufficiency Under Misordered Onboarding).} In any large-scale platform where durable identity or visibility is granted prior to the accumulation of verifiable practice, no moderation regime operating solely at the level of content review, rule enforcement, or actor removal can restore identity stabilization or semantic coherence at scale.

\textit{Sketch of argument.} Moderation operates on existing entities and their outputs. It presupposes that identities are already instantiated and that violations can be detected relative to stable categories or norms. If the platform’s onboarding architecture permits identity instantiation without practice, then the system admits an arbitrarily large population of entities whose semantic states are underdetermined by their histories. These entities generate entropy not primarily through rule violation, but through incoherence and indeterminacy. Moderation can remove or suppress individual instances, but it cannot alter the underlying rate at which new incoherent identities are created. As long as aspiration is permitted to crystallize into durable identity upstream of action, entropy injection exceeds any feasible rate of moderation-driven dissipation.

Conversely, if onboarding enforces practice before identity recognition, the entropy injection rate is bounded at the point of entry. Moderation then becomes tractable because it operates on entities whose histories already constrain interpretation. The theorem follows from a simple conservation argument: no downstream filter can compensate for an upstream generator whose output rate exceeds the system’s capacity to bind meaning.

\section{Synthesis: Tutorials, Identity Decay, and Engagement Optimization}

Taken together, the analysis of tutorial structure, identity stabilization, RSVP field dynamics, and empirical signatures converges on a single architectural diagnosis. Contemporary engagement-optimized platforms have inverted the developmental order of participation. They reward declaration before demonstration, visibility before continuity, and aspiration before practice. Tutorials, which should function as coupling mechanisms, instead operate as scalar inflators, injecting semantic potential without enforcing realization.

The result is not primarily deception, nor even failure in the moral sense. It is ontological confusion. Identity ceases to function as a binding constraint and becomes a decorative surface. Practice becomes optional. History loses its authority. In RSVP terms, scalar potential is allowed to accumulate without vector discharge, and entropy settles into a stable equilibrium rather than being dissipated through flow.

The counterexamples discussed earlier demonstrate that this outcome is not inevitable. Systems that enforce capability scaffolding, that refuse to simulate arrival, and that withhold identity until competence is demonstrated reliably produce durable skills, stable roles, and low semantic entropy. These systems succeed not because they motivate users more effectively, but because they encode the correct ordering constraints.

The broader implication is that meaning cannot be engineered through incentives alone. It emerges only when identity, action, and consequence are tightly coupled from the moment of entry. Engagement optimization that ignores this coupling will continue to generate environments that appear vibrant at the surface while hollowing out the conditions of trust beneath. Tutorials that promise arrival without enforcing becoming are not neutral affordances; they are entropy engines.

A platform that wishes to support genuine participation must therefore treat onboarding not as a marketing funnel but as a formative phase in which ontological commitments are established. Identity must be earned, not declared. Practice must precede projection. Only under these constraints can optimization operate without dissolving the very meanings it seeks to amplify.

\section{Practice-First Expression as a Stabilizing Counterexample}

A concrete illustration of capability-first participation can be found in a short video uploaded to the Facebook Page associated with \textit{1}, titled \textit{Simple Practice}. The video consists of approximately seven minutes of uninterrupted drywall work performed above a staircase. It is notable not for production polish, novelty, or aspirational framing, but for the way it foregrounds embodied competence as the sole organizing principle of the content.

The activity depicted combines several independently acquired skills into a single coherent practice. Maintaining balance on a ladder without the use of hands enables the simultaneous handling of a hawk and trowel, allowing joint compound to be applied directly rather than staged. This configuration integrates balance, fine motor control, surface reading, and iterative correction. Each pass across the surface encodes feedback from the previous pass, requiring patience, attention to edge conditions, and tolerance for gradual improvement rather than immediate completion. The video does not announce expertise; it demonstrates it. Meaning arises from visible constraint satisfaction rather than from symbolic declaration.

What is significant in the present context is not the subject matter of the video, but its structural properties. The Page hosting the video does not attempt to project a professional identity in advance of practice, nor does it frame the content as a step toward monetization or influence. The title \textit{Simple Practice} is descriptive rather than aspirational. It names the activity that is occurring rather than a status that has been achieved. As a result, the identity of the Page is implicitly stabilized by the accumulation of demonstrated capability rather than by branding rituals.

A parallel logic appears in the description of Spherepop, where instead of an explanatory narrative or marketing language, a single expression is presented: \(( (6) / ((2)((2) + (1))) )\). This representation does not explain Spherepop verbally; it performs it. The expression places an arithmetic structure into Spherepop normal form by making all scopes explicit through fully nested parentheses. The reader is not told what Spherepop is meant to be; they are shown how it operates. Understanding, if it occurs, arises through engagement with the structure rather than through identification with a projected role.

Both examples instantiate the same ordering constraint. Practice precedes identity. Scalar potential remains local and bounded because no symbolic surplus is introduced beyond what the demonstrated action can support. Vector flow is explicit and continuous, whether in the physical domain of drywalling or the symbolic domain of formal expression. Entropy is minimized not through moderation or curation, but through the refusal to grant semantic weight to anything other than realized competence.

From the RSVP perspective, these examples exhibit immediate coupling between potential and flow. There is no extended phase in which promise accumulates without discharge. As a result, there is no stable stalled state. The Page either continues to accrue meaning through additional practice or remains sparse without polluting the surrounding semantic field. Identity, such as it is, remains legible precisely because it is not prematurely asserted.

This case clarifies a critical point in the broader argument. The problem with engagement-driven platforms is not that they host novices, experimental work, or informal practice. It is that they encourage symbolic self-positioning before the conditions for meaningful practice are in place. When a system instead permits users to act, demonstrate, and iterate without requiring them to declare who they are becoming, identity stabilizes naturally as a downstream effect. The drywalling video and the Spherepop expression function not as personal anecdotes, but as minimal working examples of an architecture in which meaning is conserved by design.

\section*{Closing Synthesis}

Seen through the lens developed in this essay, the contrast between stalled aspirational Pages and practice-first expressions is no longer anecdotal but structural. The drywalling demonstration and the Spherepop normal-form expression instantiate precisely the regime in which the proposed decay law becomes unnecessary, because unsupported scalar potential is never allowed to accumulate in the first place. Where identity is not prematurely instantiated, where action is uninterrupted and legible, and where meaning is allowed to arise only as a residue of demonstrated constraint satisfaction, entropy does not require active suppression. It simply fails to stabilize. The implication is that identity decay on engagement-optimized platforms is not an unavoidable consequence of scale, nor of novice participation, but of a misordered ontology at the point of entry. When platforms reward naming before doing, potential without flow becomes metastable and decay must be imposed from above. When platforms instead allow becoming to precede declaration, stabilization emerges naturally, and meaning is conserved without governance heroics. The choice, therefore, is not between openness and coherence, but between architectures that simulate arrival and architectures that enforce practice. Only the latter can sustain identity without collapse.

\newpage
\section*{appendices}

\appendix
\section{Appendix A: Identity Stabilization as a Dynamical System}

\begin{align}
E &= \{e_i\}_{i \in \mathbb{N}} \\[6pt]
I_e(t) &\in \mathcal{N} \\[6pt]
M_e(t) &\in \mathcal{M} \\[6pt]
H_e(t) &= \bigoplus_{\tau \le t} \mathcal{P}(a_e(\tau)) \\[10pt]

M_e(t+\Delta t) &= \mathcal{T}\!\left(M_e(t), a_e(t), \Delta t ; \Pi \right) \\[10pt]

d_I\!\left(I_e(t+\Delta t), I_e(t)\right)
&\le \kappa \, \|a_e(t)\| \, \Delta t \\[10pt]

H_e(t+\Delta t) &= H_e(t) \oplus \mathcal{P}(a_e(t)) \\[10pt]

B_e(t+\Delta t) &= B_e(t) - \mathcal{C}\!\left(a_e(t)\right) \\[6pt]

d_M\!\left(M_e(t+\Delta t), M_e(t)\right)
&\le B_e(t) - B_e(t+\Delta t) \\[10pt]

\lim_{T \to \infty}
\operatorname{Var}\!\left(M_e(T)\right)
&= 0 \quad \text{iff} \quad
\int_0^T \|a_e(t)\| \, dt = \infty
\end{align}

\section{Appendix B: RSVP Field Dynamics and Identity Decay}

\begin{align}
\Phi_e(t) &\in \mathbb{R}_{\ge 0} \\[6pt]
\mathbf{v}_e(t) &\in \mathbb{R}^n \\[6pt]
S_e(t) &\in \mathbb{R}_{\ge 0} \\[10pt]

\partial_t \Phi_e
&= -\alpha \Phi_e
+ \beta \|\mathbf{v}_e\|
- \gamma \Phi_e
\,\sigma\!\left(\frac{\Phi_e}{\Phi_0}\right)
\left(1 - \sigma\!\left(\frac{\|\mathbf{v}_e\|}{v_0}\right)\right)
\\[12pt]

\partial_t \mathbf{v}_e
&= -\nabla \Phi_e
- \nu \mathbf{v}_e
+ \mathbf{F}_e(t)
\\[12pt]

\partial_t S_e
&= \eta \Phi_e
\left(1 - \sigma\!\left(\frac{\|\mathbf{v}_e\|}{v_0}\right)\right)
- \lambda \|\mathbf{v}_e\|
+ D_S \Delta S_e
\\[12pt]

\sigma(x) &= \frac{1}{1 + e^{-x}} \\[10pt]

\Phi_e(t) \gg \Phi_0 \;\wedge\; \|\mathbf{v}_e(t)\| \ll v_0
&\;\Longrightarrow\;
\partial_t \Phi_e < 0,\;
\partial_t S_e > 0
\\[12pt]

\|\mathbf{v}_e(t)\| \gg v_0
&\;\Longrightarrow\;
\partial_t \Phi_e \ge 0,\;
\partial_t S_e < 0
\\[12pt]

\lim_{t \to \infty}
\left(
\Phi_e(t), \mathbf{v}_e(t), S_e(t)
\right)
&=
\begin{cases}
(0, 0, 0), & \int_0^\infty \|\mathbf{v}_e(t)\| dt < \infty \\[6pt]
(\Phi^\ast, \mathbf{v}^\ast, S^\ast), &
\int_0^\infty \|\mathbf{v}_e(t)\| dt = \infty
\end{cases}
\end{align}

\section{Appendix C: Tutorial Ordering and Ontological Constraint}

\begin{align}
\mathcal{T}_0 &: E \rightarrow \mathcal{A} \\[6pt]
\mathcal{T}_1 &: E \rightarrow \mathcal{I} \\[10pt]

\mathcal{A}(e,t) &= \int_0^t \|a_e(\tau)\| \, d\tau \\[10pt]

\mathcal{I}(e,t) &= \mathcal{F}\!\left(H_e(t)\right) \\[10pt]

\text{Aspiration-First Ordering:} \qquad
\mathcal{T}_1 \prec \mathcal{T}_0
\\[6pt]

\text{Practice-First Ordering:} \qquad
\mathcal{T}_0 \prec \mathcal{T}_1
\\[12pt]

\mathcal{T}_1 \prec \mathcal{T}_0
&\;\Longrightarrow\;
\exists\, e :
\mathcal{I}(e,t) \neq \varnothing \;\wedge\; \mathcal{A}(e,t) \approx 0
\\[12pt]

\mathcal{T}_0 \prec \mathcal{T}_1
&\;\Longrightarrow\;
\forall\, e :
\mathcal{I}(e,t) \neq \varnothing
\;\Rightarrow\;
\mathcal{A}(e,t) \ge \epsilon
\\[12pt]

\epsilon &> 0 \\[10pt]

\partial_t \mathcal{I}(e,t)
&= \chi\!\left(\mathcal{A}(e,t) - \epsilon\right)
\\[12pt]

\chi(x) &=
\begin{cases}
0, & x < 0 \\[4pt]
1, & x \ge 0
\end{cases}
\\[12pt]

\mathcal{T}_1 \prec \mathcal{T}_0
&\;\Longrightarrow\;
\limsup_{t \to \infty}
\operatorname{Card}\!\left(
\{ e \mid \mathcal{A}(e,t) \approx 0 \}
\right)
= \infty
\\[12pt]

\mathcal{T}_0 \prec \mathcal{T}_1
&\;\Longrightarrow\;
\limsup_{t \to \infty}
\operatorname{Card}\!\left(
\{ e \mid \mathcal{A}(e,t) \approx 0 \}
\right)
< \infty
\end{align}

\section{Appendix D: Moderation Limits and Entropy Injection Rate}

\begin{align}
\Lambda(t) &= \text{rate of new identity instantiation} \\[6pt]
\mu(t) &= \text{moderation removal rate} \\[6pt]
\Sigma(t) &= \text{semantic entropy density} \\[10pt]

\dot{\Sigma}(t)
&= \int_E
\left(
\eta\,\Phi_e(t)\,\mathbb{I}\{\|\mathbf{v}_e(t)\| < v_0\}
- \lambda\,\|\mathbf{v}_e(t)\|
\right)
\,de
- \mu(t)
\\[12pt]

\Lambda(t)
&= \int_E
\delta\!\left(
\mathcal{I}(e,t^+) \neq \varnothing \;\wedge\; \mathcal{A}(e,t) \approx 0
\right)
\,de
\\[12pt]

\mu(t)
&\le \mu_{\max}
\\[10pt]

\exists\, c > 0 :
\Lambda(t) > \mu_{\max} + c
&\;\Longrightarrow\;
\dot{\Sigma}(t) > 0
\\[12pt]

\limsup_{t \to \infty} \Sigma(t)
&=
\begin{cases}
\infty, & \Lambda(t) > \mu_{\max} \\[6pt]
\Sigma^\ast < \infty, & \Lambda(t) \le \mu_{\max}
\end{cases}
\\[14pt]

\mathcal{T}_1 \prec \mathcal{T}_0
&\;\Longrightarrow\;
\Lambda(t) = \Omega(|E|)
\\[10pt]

\mathcal{T}_0 \prec \mathcal{T}_1
&\;\Longrightarrow\;
\Lambda(t) = O(1)
\\[12pt]

\forall\, \mathcal{M} :
\quad
\exists\, \Lambda^\ast :
\Lambda(t) > \Lambda^\ast
&\;\Rightarrow\;
\text{moderation failure}
\end{align}

\section{Appendix E: Provenance, Drift Budgets, and Namespace Geometry}

\begin{align}
\mathcal{N} &= \text{namespace space} \\[6pt]
\mathcal{G} &= (\mathcal{N}, \sim) \quad \text{groupoid of identities} \\[10pt]

e_1 \sim e_2
&\;\Longleftrightarrow\;
H_{e_1}(t) = H_{e_2}(t)
\\[12pt]

\pi : \mathcal{G} \rightarrow \mathcal{N}
\\[10pt]

\gamma_e : [0,t] \rightarrow \mathcal{N}
\\[6pt]

\gamma_e(t) &= \pi\!\left(H_e(t)\right)
\\[12pt]

\ell(\gamma_e)
&= \int_0^t
\left\|
\frac{d}{d\tau} \gamma_e(\tau)
\right\|
\, d\tau
\\[12pt]

B_e(t)
&= B_e(0) - \ell(\gamma_e)
\\[12pt]

B_e(t) \ge 0
\\[10pt]

B_e(t) = 0
&\;\Longrightarrow\;
\frac{d}{dt}\gamma_e(t) = 0
\\[12pt]

d_{\mathcal{N}}\!\left(
\gamma_e(t_1), \gamma_e(t_2)
\right)
&\le
\int_{t_1}^{t_2}
\left\|
\frac{d}{dt}\gamma_e(t)
\right\|
dt
\\[14pt]

\exists\, L < \infty :
\forall\, e,\;
\ell(\gamma_e) \le L
\\[12pt]

\lim_{t \to \infty}
d_{\mathcal{N}}\!\left(
\gamma_e(t), \gamma_e(t-\Delta t)
\right)
&= 0
\quad \text{iff} \quad
\int_0^\infty \|a_e(t)\| dt = \infty
\\[14pt]

\mathcal{N}_{\text{stable}}
&=
\left\{
e \in E \mid
\lim_{t \to \infty}
\frac{d}{dt}\gamma_e(t) = 0
\right\}
\\[12pt]

\operatorname{Card}(\mathcal{N}_{\text{stable}})
&< \infty
\quad \text{under practice-first ordering}
\\[10pt]

\operatorname{Card}(\mathcal{N}_{\text{stable}})
&= \infty
\quad \text{under aspiration-first ordering}
\end{align}

\section{Appendix F: Discovery Dynamics and Signal-to-Noise Collapse}

\begin{align}
\mathcal{D}(e,t) &= \text{discovery probability of entity } e \\[6pt]
\mathcal{R}(t) &= \text{recommendation operator} \\[6pt]
\mathcal{S}(e,t) &= \text{semantic signal strength} \\[6pt]
\mathcal{N}(t) &= \text{semantic noise floor} \\[10pt]

\mathcal{D}(e,t)
&= \mathcal{R}(t)\!\left(
\mathcal{S}(e,t) - \mathcal{N}(t)
\right)
\\[12pt]

\mathcal{S}(e,t)
&= \int_0^t
w(\tau)\,\|a_e(\tau)\|\,d\tau
\\[12pt]

\mathcal{N}(t)
&= \sum_{e \in E}
\mathbb{I}\{\|\mathbf{v}_e(t)\| < v_0\}
\\[12pt]

\mathrm{SNR}(t)
&=
\frac{
\sum_{e \in E}
\mathcal{S}(e,t)
}{
\mathcal{N}(t)
}
\\[12pt]

\partial_t \mathcal{N}(t)
&= \Lambda(t) - \mu(t)
\\[10pt]

\Lambda(t) > \mu(t)
&\;\Longrightarrow\;
\partial_t \mathrm{SNR}(t) < 0
\\[12pt]

\lim_{t \to \infty}
\mathrm{SNR}(t)
&=
\begin{cases}
0, & \Lambda(t) > \mu_{\max} \\[6pt]
\mathrm{SNR}^\ast > 0, & \Lambda(t) \le \mu_{\max}
\end{cases}
\\[14pt]

\mathcal{T}_1 \prec \mathcal{T}_0
&\;\Longrightarrow\;
\mathcal{N}(t) = \Omega(|E|)
\\[10pt]

\mathcal{T}_0 \prec \mathcal{T}_1
&\;\Longrightarrow\;
\mathcal{N}(t) = O(1)
\\[12pt]

\forall\, e :
\quad
\lim_{t \to \infty}
\mathcal{D}(e,t) = 0
\;\;\text{iff}\;\;
\mathrm{SNR}(t) = 0
\end{align}

\section{Appendix G: Ontological Phase Transitions and Platform Regimes}

\begin{align}
\Theta(t) &= \text{global onboarding ordering parameter} \\[6pt]
\Theta(t) &\in [0,1] \\[10pt]

\Theta(t) = 0
&\;\Longleftrightarrow\;
\mathcal{T}_0 \prec \mathcal{T}_1
\\[6pt]

\Theta(t) = 1
&\;\Longleftrightarrow\;
\mathcal{T}_1 \prec \mathcal{T}_0
\\[12pt]

\mathcal{Z}(t)
&=
\sum_{e \in E}
\exp\!\left(
-\frac{\|\mathbf{v}_e(t)\|}{v_0}
\right)
\\[12pt]

\mathcal{F}(\Theta)
&=
\lim_{t \to \infty}
\frac{1}{|E|}
\mathcal{Z}(t)
\\[12pt]

\exists\, \Theta_c \in (0,1) :
\quad
\lim_{\epsilon \to 0^+}
\mathcal{F}(\Theta_c + \epsilon)
\neq
\lim_{\epsilon \to 0^+}
\mathcal{F}(\Theta_c - \epsilon)
\\[14pt]

\Theta < \Theta_c
&\;\Longrightarrow\;
\limsup_{t \to \infty}
\operatorname{Card}(\mathcal{N}_{\text{stable}})
< \infty
\\[10pt]

\Theta > \Theta_c
&\;\Longrightarrow\;
\limsup_{t \to \infty}
\operatorname{Card}(\mathcal{N}_{\text{stable}})
= \infty
\\[14pt]

\partial_\Theta \mathcal{F}(\Theta)
&=
\begin{cases}
< 0, & \Theta < \Theta_c \\[6pt]
= \infty, & \Theta = \Theta_c \\[6pt]
> 0, & \Theta > \Theta_c
\end{cases}
\\[14pt]

\Sigma(t)
&\sim
\begin{cases}
O(1), & \Theta < \Theta_c \\[6pt]
O(|E|), & \Theta > \Theta_c
\end{cases}
\\[14pt]

\Theta(t) \searrow \Theta_c
&\;\Longrightarrow\;
\text{critical slowing down}
\\[10pt]

\Theta(t) \nearrow \Theta_c
&\;\Longrightarrow\;
\text{entropy runaway}
\end{align}

\section{Appendix H: Semantic Renormalization and Scale Separation}

\begin{align}
\mathcal{E}_\ell &= \{ e \in E \mid \|a_e(t)\| \sim \ell \} \\[6pt]
\ell &\in \mathbb{R}_{>0} \\[10pt]

\Phi_\ell(t)
&=
\int_{\mathcal{E}_\ell}
\Phi_e(t)\,de
\\[12pt]

\mathbf{v}_\ell(t)
&=
\int_{\mathcal{E}_\ell}
\mathbf{v}_e(t)\,de
\\[12pt]

S_\ell(t)
&=
\int_{\mathcal{E}_\ell}
S_e(t)\,de
\\[12pt]

\partial_t \Phi_\ell
&=
-\alpha \Phi_\ell
+ \beta \|\mathbf{v}_\ell\|
- \gamma \Phi_\ell
\left(
1 - \sigma\!\left(
\frac{\|\mathbf{v}_\ell\|}{v_0}
\right)
\right)
\\[12pt]

\partial_t S_\ell
&=
\eta \Phi_\ell
\left(
1 - \sigma\!\left(
\frac{\|\mathbf{v}_\ell\|}{v_0}
\right)
\right)
- \lambda \|\mathbf{v}_\ell\|
+ D_\ell \Delta S_\ell
\\[12pt]

\ell_1 < \ell_2
&\;\Longrightarrow\;
\tau_{\ell_1} \ll \tau_{\ell_2}
\\[10pt]

\tau_\ell
&=
\left(
\alpha +
\gamma\left(
1 - \sigma\!\left(
\frac{\|\mathbf{v}_\ell\|}{v_0}
\right)
\right)
\right)^{-1}
\\[14pt]

\exists\, \ell^\ast :
\forall\, \ell < \ell^\ast,
\quad
\lim_{t \to \infty} \Phi_\ell(t) = 0
\\[10pt]

\exists\, \ell^\ast :
\forall\, \ell > \ell^\ast,
\quad
\lim_{t \to \infty} \Phi_\ell(t) = \Phi_\ell^\ast > 0
\\[14pt]

\Phi_{\text{eff}}
&=
\int_{\ell^\ast}^{\infty}
\Phi_\ell \, d\ell
\\[12pt]

S_{\text{eff}}
&=
\int_{\ell^\ast}^{\infty}
S_\ell \, d\ell
\\[12pt]

\frac{d}{dt} S_{\text{eff}} < 0
&\;\Longleftrightarrow\;
\Theta < \Theta_c
\\[10pt]

\frac{d}{dt} S_{\text{eff}} > 0
&\;\Longleftrightarrow\;
\Theta > \Theta_c
\end{align}

\section{Appendix I: Fixed Points, Attractors, and Collapse Modes}

\begin{align}
\mathbf{X}_e(t)
&=
\begin{pmatrix}
\Phi_e(t) \\
\mathbf{v}_e(t) \\
S_e(t)
\end{pmatrix}
\\[10pt]

\dot{\mathbf{X}}_e
&=
\mathcal{F}\!\left(\mathbf{X}_e; \Theta\right)
\\[12pt]

\mathbf{X}_e^\ast
&:\;
\mathcal{F}\!\left(\mathbf{X}_e^\ast; \Theta\right) = 0
\\[12pt]

\mathbf{X}_0^\ast
&=
\begin{pmatrix}
0 \\
\mathbf{0} \\
0
\end{pmatrix}
\\[12pt]

\mathbf{X}_1^\ast
&=
\begin{pmatrix}
\Phi^\ast \\
\mathbf{v}^\ast \\
S^\ast
\end{pmatrix}
\quad
\text{with }
\|\mathbf{v}^\ast\| > v_0
\\[14pt]

J(\mathbf{X})
&=
\frac{\partial \mathcal{F}}{\partial \mathbf{X}}
\\[10pt]

\operatorname{spec}
\left(
J(\mathbf{X}_0^\ast)
\right)
&=
\begin{cases}
\Re(\lambda_i) < 0, & \Theta < \Theta_c \\[6pt]
\exists\, \lambda_i > 0, & \Theta > \Theta_c
\end{cases}
\\[14pt]

\operatorname{spec}
\left(
J(\mathbf{X}_1^\ast)
\right)
&\subset
\{ \lambda \in \mathbb{C} \mid \Re(\lambda) < 0 \}
\\[14pt]

\mathcal{A}_0
&=
\left\{
\mathbf{X} \mid
\lim_{t \to \infty}
\mathbf{X}(t) = \mathbf{X}_0^\ast
\right\}
\\[12pt]

\mathcal{A}_1
&=
\left\{
\mathbf{X} \mid
\lim_{t \to \infty}
\mathbf{X}(t) = \mathbf{X}_1^\ast
\right\}
\\[14pt]

\Theta < \Theta_c
&\;\Longrightarrow\;
\mu(\mathcal{A}_1) > 0
\\[10pt]

\Theta > \Theta_c
&\;\Longrightarrow\;
\mu(\mathcal{A}_1) = 0
\\[14pt]

\Theta > \Theta_c
&\;\Longrightarrow\;
\exists\, \mathcal{A}_\infty :
\quad
\lim_{t \to \infty} S_e(t) = \infty
\\[14pt]

\mathcal{A}_\infty
&=
\left\{
\mathbf{X} \mid
\lim_{t \to \infty}
\Phi_e(t) > 0,\;
\lim_{t \to \infty}
\|\mathbf{v}_e(t)\| = 0
\right\}
\\[14pt]

\Theta \searrow \Theta_c
&\;\Longrightarrow\;
\text{basin fracturing}
\\[10pt]

\Theta \nearrow \Theta_c
&\;\Longrightarrow\;
\text{attractor annihilation}
\end{align}

\section{Appendix I: Fixed Points, Attractors, and Collapse Modes}

\begin{align}
\mathbf{X}_e(t)
&=
\begin{pmatrix}
\Phi_e(t) \\
\mathbf{v}_e(t) \\
S_e(t)
\end{pmatrix}
\\[10pt]

\dot{\mathbf{X}}_e
&=
\mathcal{F}\!\left(\mathbf{X}_e; \Theta\right)
\\[12pt]

\mathbf{X}_e^\ast
&:\;
\mathcal{F}\!\left(\mathbf{X}_e^\ast; \Theta\right) = 0
\\[12pt]

\mathbf{X}_0^\ast
&=
\begin{pmatrix}
0 \\
\mathbf{0} \\
0
\end{pmatrix}
\\[12pt]

\mathbf{X}_1^\ast
&=
\begin{pmatrix}
\Phi^\ast \\
\mathbf{v}^\ast \\
S^\ast
\end{pmatrix}
\quad
\text{with }
\|\mathbf{v}^\ast\| > v_0
\\[14pt]

J(\mathbf{X})
&=
\frac{\partial \mathcal{F}}{\partial \mathbf{X}}
\\[10pt]

\operatorname{spec}
\left(
J(\mathbf{X}_0^\ast)
\right)
&=
\begin{cases}
\Re(\lambda_i) < 0, & \Theta < \Theta_c \\[6pt]
\exists\, \lambda_i > 0, & \Theta > \Theta_c
\end{cases}
\\[14pt]

\operatorname{spec}
\left(
J(\mathbf{X}_1^\ast)
\right)
&\subset
\{ \lambda \in \mathbb{C} \mid \Re(\lambda) < 0 \}
\\[14pt]

\mathcal{A}_0
&=
\left\{
\mathbf{X} \mid
\lim_{t \to \infty}
\mathbf{X}(t) = \mathbf{X}_0^\ast
\right\}
\\[12pt]

\mathcal{A}_1
&=
\left\{
\mathbf{X} \mid
\lim_{t \to \infty}
\mathbf{X}(t) = \mathbf{X}_1^\ast
\right\}
\\[14pt]

\Theta < \Theta_c
&\;\Longrightarrow\;
\mu(\mathcal{A}_1) > 0
\\[10pt]

\Theta > \Theta_c
&\;\Longrightarrow\;
\mu(\mathcal{A}_1) = 0
\\[14pt]

\Theta > \Theta_c
&\;\Longrightarrow\;
\exists\, \mathcal{A}_\infty :
\quad
\lim_{t \to \infty} S_e(t) = \infty
\\[14pt]

\mathcal{A}_\infty
&=
\left\{
\mathbf{X} \mid
\lim_{t \to \infty}
\Phi_e(t) > 0,\;
\lim_{t \to \infty}
\|\mathbf{v}_e(t)\| = 0
\right\}
\\[14pt]

\Theta \searrow \Theta_c
&\;\Longrightarrow\;
\text{basin fracturing}
\\[10pt]

\Theta \nearrow \Theta_c
&\;\Longrightarrow\;
\text{attractor annihilation}
\end{align}

\section{Appendix K: Derivation of the Potential Decay Law from a Lyapunov Functional}

\begin{align}
\mathbf{X}_e
&=
\begin{pmatrix}
\Phi_e \\
\mathbf{v}_e \\
S_e
\end{pmatrix}
\\[10pt]

\mathcal{E}_e(\Phi_e,\mathbf{v}_e,S_e)
&=
\frac{1}{2}\|\mathbf{v}_e\|^2
+ U(\Phi_e)
+ W(S_e)
+ C(\Phi_e,\mathbf{v}_e)
\\[12pt]

U(\Phi)
&=
\frac{\alpha}{2}\Phi^2
\\[10pt]

W(S)
&=
\frac{\omega}{2}S^2
\\[10pt]

C(\Phi,\mathbf{v})
&=
\gamma\,\Phi^2
\left(1-\sigma\!\left(\frac{\|\mathbf{v}\|}{v_0}\right)\right)
\\[12pt]

\dot{\Phi}_e
&=
-\partial_{\Phi}\mathcal{E}_e
+
\beta \|\mathbf{v}_e\|
\\[10pt]

\dot{\mathbf{v}}_e
&=
-\partial_{\mathbf{v}}\mathcal{E}_e
\\[10pt]

\dot{S}_e
&=
-\partial_{S}\mathcal{E}_e
+
\eta\,\Phi_e
\left(1-\sigma\!\left(\frac{\|\mathbf{v}_e\|}{v_0}\right)\right)
\\[14pt]

\partial_{\Phi}\mathcal{E}_e
&=
\alpha\Phi_e
+
2\gamma\Phi_e
\left(1-\sigma\!\left(\frac{\|\mathbf{v}_e\|}{v_0}\right)\right)
\\[12pt]

\Rightarrow\quad
\dot{\Phi}_e
&=
-\alpha\Phi_e
-
2\gamma\Phi_e
\left(1-\sigma\!\left(\frac{\|\mathbf{v}_e\|}{v_0}\right)\right)
+
\beta \|\mathbf{v}_e\|
\\[14pt]

\partial_{\mathbf{v}}\mathcal{E}_e
&=
\mathbf{v}_e
+
\gamma\Phi_e^2
\left(-\sigma'\!\left(\frac{\|\mathbf{v}_e\|}{v_0}\right)\right)
\frac{1}{v_0}\frac{\mathbf{v}_e}{\|\mathbf{v}_e\|}
\\[12pt]

\Rightarrow\quad
\dot{\mathbf{v}}_e
&=
-\mathbf{v}_e
+
\frac{\gamma\Phi_e^2}{v_0}
\sigma'\!\left(\frac{\|\mathbf{v}_e\|}{v_0}\right)
\frac{\mathbf{v}_e}{\|\mathbf{v}_e\|}
\\[14pt]

\partial_S \mathcal{E}_e
&=
\omega S_e
\\[10pt]

\Rightarrow\quad
\dot{S}_e
&=
-\omega S_e
+
\eta\,\Phi_e
\left(1-\sigma\!\left(\frac{\|\mathbf{v}_e\|}{v_0}\right)\right)
\\[16pt]

\frac{d}{dt}\mathcal{E}_e
&=
\partial_{\Phi}\mathcal{E}_e\,\dot{\Phi}_e
+
\left\langle
\partial_{\mathbf{v}}\mathcal{E}_e,\dot{\mathbf{v}}_e
\right\rangle
+
\partial_S\mathcal{E}_e\,\dot{S}_e
\\[12pt]

&=
-\left(\partial_{\Phi}\mathcal{E}_e\right)^2
-\left\|
\partial_{\mathbf{v}}\mathcal{E}_e
\right\|^2
-\left(\partial_S\mathcal{E}_e\right)^2
+
\beta \|\mathbf{v}_e\|\,\partial_{\Phi}\mathcal{E}_e
+
\eta\,\Phi_e
\left(1-\sigma\!\left(\frac{\|\mathbf{v}_e\|}{v_0}\right)\right)
\partial_S\mathcal{E}_e
\\[14pt]

\beta \|\mathbf{v}_e\|\,\partial_{\Phi}\mathcal{E}_e
&\le
\frac{1}{2}\left(\partial_{\Phi}\mathcal{E}_e\right)^2
+\frac{\beta^2}{2}\|\mathbf{v}_e\|^2
\\[12pt]

\eta\,\Phi_e
\left(1-\sigma\!\left(\frac{\|\mathbf{v}_e\|}{v_0}\right)\right)
\partial_S\mathcal{E}_e
&\le
\frac{1}{2}\left(\partial_S\mathcal{E}_e\right)^2
+\frac{\eta^2}{2}\Phi_e^2
\left(1-\sigma\!\left(\frac{\|\mathbf{v}_e\|}{v_0}\right)\right)^2
\\[14pt]

\Rightarrow\quad
\frac{d}{dt}\mathcal{E}_e
&\le
-\frac{1}{2}\left(\partial_{\Phi}\mathcal{E}_e\right)^2
-\left\|
\partial_{\mathbf{v}}\mathcal{E}_e
\right\|^2
-\frac{1}{2}\left(\partial_S\mathcal{E}_e\right)^2
+
\frac{\beta^2}{2}\|\mathbf{v}_e\|^2
+
\frac{\eta^2}{2}\Phi_e^2
\left(1-\sigma\!\left(\frac{\|\mathbf{v}_e\|}{v_0}\right)\right)^2
\end{align}

\section{Appendix L: Proof of Monotone Entropy Growth in the Stalled Regime}

\begin{align}
\dot{S}_e
&=
-\omega S_e
+
\eta\,\Phi_e
\left(1-\sigma\!\left(\frac{\|\mathbf{v}_e\|}{v_0}\right)\right)
\\[12pt]

\text{Assume}\quad
\|\mathbf{v}_e(t)\| \le v_0/2
&\;\Rightarrow\;
\sigma\!\left(\frac{\|\mathbf{v}_e(t)\|}{v_0}\right)
\le
\sigma\!\left(\frac{1}{2}\right)
\\[10pt]

c_0
&=
1-\sigma\!\left(\frac{1}{2}\right)
> 0
\\[12pt]

\Rightarrow\quad
\dot{S}_e
&\ge
-\omega S_e
+
\eta\,c_0\,\Phi_e
\\[14pt]

\dot{\Phi}_e
&=
-\alpha\Phi_e
-
2\gamma\Phi_e
\left(1-\sigma\!\left(\frac{\|\mathbf{v}_e\|}{v_0}\right)\right)
+
\beta \|\mathbf{v}_e\|
\\[12pt]

\|\mathbf{v}_e(t)\| \le v_0/2
&\Rightarrow
1-\sigma\!\left(\frac{\|\mathbf{v}_e(t)\|}{v_0}\right)
\ge
c_0
\\[10pt]

\Rightarrow\quad
\dot{\Phi}_e
&\le
-\left(\alpha+2\gamma c_0\right)\Phi_e
+
\beta \frac{v_0}{2}
\\[14pt]

k
&=
\alpha+2\gamma c_0
\\[10pt]

\Rightarrow\quad
\Phi_e(t)
&\le
\Phi_e(0)e^{-kt}
+
\frac{\beta v_0}{2k}\left(1-e^{-kt}\right)
\\[16pt]

\dot{S}_e
&\ge
-\omega S_e
+
\eta c_0
\left[
\Phi_e(0)e^{-kt}
+
\frac{\beta v_0}{2k}\left(1-e^{-kt}\right)
\right]
\\[14pt]

S_e(t)
&\ge
S_e(0)e^{-\omega t}
+
\eta c_0
\int_0^t
e^{-\omega(t-\tau)}
\left[
\Phi_e(0)e^{-k\tau}
+
\frac{\beta v_0}{2k}\left(1-e^{-k\tau}\right)
\right]d\tau
\\[16pt]

\liminf_{t\to\infty} S_e(t)
&\ge
\eta c_0
\left(
\frac{\beta v_0}{2k}
\right)
\frac{1}{\omega}
\\[14pt]

S_e(0)=0,\;\Phi_e(0)>0,\;\beta=0
&\Rightarrow
S_e(t)
=
\eta c_0\Phi_e(0)\frac{e^{-kt}-e^{-\omega t}}{\omega-k}
\quad(\omega\ne k)
\\[14pt]

\omega < k
&\Rightarrow
\exists\, t_1:\ \forall t>t_1,\ \dot{S}_e(t)>0
\end{align}

\section{Appendix M: Proof of Moderation Insufficiency via a Conservation Bound}

\begin{align}
\mathcal{N}(t)
&=
\operatorname{Card}\!\left(
\left\{
e \in E:
\mathcal{I}(e,t)\neq\varnothing\ \wedge\ \mathcal{A}(e,t)\le \epsilon
\right\}
\right)
\\[12pt]

\dot{\mathcal{N}}(t)
&=
\Lambda(t) - \mu(t)
\\[12pt]

\mu(t)
&\le
\mu_{\max}
\\[12pt]

\Lambda(t)
&\ge
\Lambda_{\min}
> \mu_{\max}
\\[12pt]

\Rightarrow\quad
\dot{\mathcal{N}}(t)
&\ge
\Lambda_{\min} - \mu_{\max}
= c
> 0
\\[14pt]

\Rightarrow\quad
\mathcal{N}(t)
&\ge
\mathcal{N}(0) + ct
\\[14pt]

\mathcal{N}(t) \to \infty
&\quad\text{as}\quad t\to\infty
\\[16pt]

\Sigma(t)
&=
\int_E
S_e(t)\,de
\\[12pt]

S_e(t)
&\ge s_0\,
\mathbb{I}\!\left\{
\mathcal{I}(e,t)\neq\varnothing\ \wedge\ \mathcal{A}(e,t)\le \epsilon
\right\}
\qquad (s_0>0)
\\[12pt]

\Rightarrow\quad
\Sigma(t)
&\ge
s_0\,\mathcal{N}(t)
\\[12pt]

\Rightarrow\quad
\Sigma(t)
&\ge
s_0\mathcal{N}(0) + s_0 c t
\\[12pt]

\Sigma(t) \to \infty
&\quad\text{as}\quad t\to\infty
\\[16pt]

\forall\, \mu(\cdot)\le \mu_{\max},\;
\exists\, \Lambda_{\min}>\mu_{\max}
&\Rightarrow
\lim_{t\to\infty}\Sigma(t)=\infty
\end{align}

\section{Appendix N: Groupoids, Fibrations, and Provenance as a Functor}

\begin{align}
\mathsf{E} &: \text{small category of entities} \\[4pt]
\mathrm{Ob}(\mathsf{E}) &= E \\[4pt]
\mathrm{Hom}_{\mathsf{E}}(e,e') &= \{\text{typed transformations } e \to e'\} \\[10pt]

\mathsf{H} &: \text{free (path) category on transformation generators} \\[4pt]
\mathrm{Ob}(\mathsf{H}) &= E \\[4pt]
\mathrm{Hom}_{\mathsf{H}}(e,e') &= \{\text{finite composable sequences } a_1;\dots;a_k\} \\[10pt]

Q &: \mathsf{H} \to \mathsf{E} \qquad \text{(quotient by relations)} \\[10pt]

\Pi &: \mathsf{E} \to \mathsf{N} \qquad \text{(namespace functor)} \\[6pt]
\mathsf{N} &: \text{groupoid of namespaces} \\[10pt]

F &= \Pi \circ Q : \mathsf{H} \to \mathsf{N} \\[14pt]

\mathsf{G} &= \mathrm{Core}(\mathsf{E}) \qquad \text{(maximal subgroupoid)} \\[10pt]

p &: \mathsf{G} \to \mathsf{N} \qquad \text{(groupoid fibration)} \\[12pt]

\forall\, u:n\to n' \in \mathrm{Hom}_{\mathsf{N}},\ \forall\, e \in p^{-1}(n),\ 
\exists\, \tilde{u}:e\to e' \in \mathrm{Hom}_{\mathsf{G}}
&\ \text{s.t.}\ p(\tilde{u})=u
\\[14pt]

\mathsf{A} &: \text{action monoid as one-object category} \\[4pt]
\mathrm{Ob}(\mathsf{A}) &= \{\ast\} \\[4pt]
\mathrm{End}_{\mathsf{A}}(\ast) &= (A,\cdot,1) \\[10pt]

\alpha_e &: \mathsf{A} \to \mathsf{H} \qquad \text{(action trace functor for entity $e$)} \\[10pt]

H_e &= \alpha_e(A) \in \mathrm{End}_{\mathsf{H}}(e) \\[12pt]

I_e &= F(H_e) \in \mathrm{End}_{\mathsf{N}}(n_e) \\[16pt]

\mathsf{V} &: \text{poset-enriched category of budgets} \\[6pt]
\mathrm{Ob}(\mathsf{V}) &= \mathbb{R}_{\ge 0} \\[4pt]
x \to y &\Longleftrightarrow x \ge y \\[10pt]

\ell &: \mathrm{Mor}(\mathsf{H}) \to \mathbb{R}_{\ge 0} \qquad \text{(cost/length)} \\[10pt]

\ell(g\circ f) &= \ell(g)+\ell(f) \\[10pt]

B_e(t)
&=
B_0 - \ell\!\left(H_e(t)\right)
\\[12pt]

B_e(t) \ge 0
&\Longleftrightarrow
\ell\!\left(H_e(t)\right) \le B_0
\\[14pt]

d_{\mathsf{N}}(n,n')
&=
\inf\left\{
\ell(f) \mid f\in \mathrm{Hom}_{\mathsf{H}}(e,e'),\ \Pi(Q(f))(n)=n'
\right\}
\\[14pt]

\mathrm{Aut}_{\mathsf{N}}(n)
&=
\{u:n\to n\}
\\[12pt]

\mathrm{Stab}(n)
&=
\left\{
e\in E \mid p(e)=n,\ \exists\, t:\ \ell(H_e(t+\Delta t))-\ell(H_e(t))\to 0
\right\}
\end{align}

\section{Appendix O: Sheaves of Meaning, Gluing, and Obstruction to Coherence}

\begin{align}
X &: \text{semantic space (site)} \\[4pt]
\mathcal{O}(X) &: \text{category of opens with inclusions} \\[10pt]

\mathcal{C} &: \text{category of semantic states} \\[10pt]

\mathcal{F} &: \mathcal{O}(X)^{\mathrm{op}} \to \mathcal{C} \qquad \text{(presheaf of meanings)} \\[12pt]

\rho_{UV} &: \mathcal{F}(U) \to \mathcal{F}(V) \qquad (V\subseteq U) \\[10pt]

\rho_{UU} &= \mathrm{id} \\[6pt]
\rho_{VW}\circ \rho_{UV} &= \rho_{UW} \qquad (W\subseteq V\subseteq U) \\[14pt]

\{U_i\}_{i\in I} &: \text{cover of } U \\[10pt]

s_i \in \mathcal{F}(U_i),\ \rho_{U_i,U_i\cap U_j}(s_i)=\rho_{U_j,U_i\cap U_j}(s_j)
&\Longrightarrow
\exists!\ s\in \mathcal{F}(U)\ \text{s.t.}\ \rho_{U,U_i}(s)=s_i
\\[14pt]

\mathrm{Eq}\!\left(
\prod_i \mathcal{F}(U_i)
\rightrightarrows
\prod_{i,j}\mathcal{F}(U_i\cap U_j)
\right)
&\cong
\mathcal{F}(U)
\\[16pt]

\mathcal{F}_e &: \mathcal{O}(X)^{\mathrm{op}} \to \mathcal{C} \qquad \text{(entity-indexed presheaf)} \\[10pt]

\mathcal{F}_e(U)
&=
\{ \text{semantic assertions about $e$ valid on context $U$} \}
\\[12pt]

\delta_{ij}
&=
\rho_{U_i,U_i\cap U_j}(s_i)\cdot
\rho_{U_j,U_i\cap U_j}(s_j)^{-1}
\in
\mathrm{Aut}\big(\mathcal{F}_e(U_i\cap U_j)\big)
\\[14pt]

\delta_{ij}\delta_{jk}\delta_{ki}
&=
1
\qquad \text{on } U_i\cap U_j\cap U_k
\\[14pt]

[\delta]
&\in
H^1\!\left(U,\underline{\mathrm{Aut}}(\mathcal{F}_e)\right)
\\[14pt]

[\delta]=0
&\Longleftrightarrow
\exists\ \text{global section } s\in \mathcal{F}_e(U)
\\[16pt]

\mathcal{P}_e &: \text{provenance sheaf} \\[6pt]
\mathcal{P}_e(U) &=
\{\text{typed transformation histories over } U\}
\\[12pt]

\mathcal{P}_e \ \text{sheaf}
&\Longleftrightarrow
\mathrm{Eq}\!\left(
\prod_i \mathcal{P}_e(U_i)
\rightrightarrows
\prod_{i,j}\mathcal{P}_e(U_i\cap U_j)
\right)
\cong
\mathcal{P}_e(U)
\\[16pt]

\mathcal{F}_e
&=
\mathsf{Sem}\circ \mathcal{P}_e
\\[10pt]

\mathsf{Sem} &: \text{functor from provenance to semantics} \\[14pt]

\text{Coherence}(e,U)
&\Longleftrightarrow
[\delta]=0
\\[14pt]

\text{Drift}(e,U)
&\Longleftrightarrow
[\delta]\neq 0
\\[16pt]

\mathcal{F}_e \ \text{stack}
&\Longleftrightarrow
\text{effective descent for objects and morphisms}
\\[10pt]

\mathrm{Desc}\!\left(\{U_i\}\right)
&=
\left(
\{s_i\},\{\phi_{ij}\},\phi_{ij}\phi_{jk}=\phi_{ik}
\right)
\\[12pt]

\mathrm{Eff}\!\left(\mathrm{Desc}(\{U_i\})\right)
&\cong
\mathcal{F}_e(U)
\end{align}

\section{Appendix P: Drift Budget as a Sheaf Metric and Renaming Cocycle}

\begin{align}
d_U &: \mathcal{F}_e(U)\times \mathcal{F}_e(U) \to \mathbb{R}_{\ge 0}
\\[10pt]

d_{U_i\cap U_j}\!\left(
\rho_{U_i,U_i\cap U_j}(s_i),
\rho_{U_j,U_i\cap U_j}(s_j)
\right)
&=
\ell(\delta_{ij})
\\[12pt]

\ell(\delta_{ij})
&=
\ell(\delta_{ji})
\\[10pt]

\ell(\delta_{ij}\delta_{jk})
&\le
\ell(\delta_{ij})+\ell(\delta_{jk})
\\[14pt]

B_e(U)
&=
B_0 - \sum_{(i,j)} \ell(\delta_{ij})
\\[12pt]

B_e(U)\ge 0
&\iff
\sum_{(i,j)} \ell(\delta_{ij}) \le B_0
\\[14pt]

\delta_{ij}=1\ \forall i,j
&\iff
[\delta]=0
\\[12pt]

[\delta]\neq 0
&\Rightarrow
\text{obstruction to gluing}
\\[12pt]

[\delta]\neq 0 \;\wedge\; B_e(U)=0
&\Rightarrow
\Phi_e \to 0
\end{align}

\section{Appendix Q: Spherepop Calculus --- Concrete Syntax (EBNF)}

\[
\begin{array}{rcl}

\langle \text{Expr} \rangle
&::=&
\langle \text{Add} \rangle
\\[6pt]

\langle \text{Add} \rangle
&::=&
\langle \text{Mul} \rangle
\;\{\; (\texttt{+}\mid\texttt{-})\; \langle \text{Mul} \rangle \;\}
\\[6pt]

\langle \text{Mul} \rangle
&::=&
\langle \text{Pow} \rangle
\;\{\; (\texttt{*}\mid\texttt{/})\; \langle \text{Pow} \rangle \;\}
\\[6pt]

\langle \text{Pow} \rangle
&::=&
\langle \text{App} \rangle
\;\{\; \texttt{^}\; \langle \text{Pow} \rangle \;\}
\\[6pt]

\langle \text{App} \rangle
&::=&
\langle \text{Atom} \rangle
\;\{\; \langle \text{Atom} \rangle \;\}
\\[6pt]

\langle \text{Atom} \rangle
&::=&
\langle \text{Nat} \rangle
\mid
\langle \text{Var} \rangle
\mid
\texttt{(}\,\langle \text{Expr} \rangle\,\texttt{)}
\\[6pt]

\langle \text{Nat} \rangle
&::=&
\texttt{0}
\mid
(\texttt{1}\mid\cdots\mid\texttt{9})\{\texttt{0}\mid\cdots\mid\texttt{9}\}
\\[6pt]

\langle \text{Var} \rangle
&::=&
(\texttt{a}\mid\cdots\mid\texttt{z}\mid\texttt{A}\mid\cdots\mid\texttt{Z})
\{\texttt{a}\mid\cdots\mid\texttt{z}\mid\texttt{A}\mid\cdots\mid\texttt{Z}\mid\texttt{0}\mid\cdots\mid\texttt{9}\mid\texttt{\_}\}

\end{array}
\]


\section{Appendix R: Spherepop Calculus --- Abstract Syntax and Scopes}

\begin{align}
t &\in \mathsf{Term} \\[6pt]
t &::= \mathsf{Num}(n)\mid \mathsf{Var}(x)\mid \mathsf{App}(t,t)\mid \mathsf{Bin}(o,t,t)\mid \mathsf{Par}(t) \\[8pt]
n &\in \mathbb{N},\quad x\in \mathsf{VarName},\quad o\in \{+,-,\times,\div,\wedge\} \\[12pt]

\operatorname{scope}(t) &\in \mathsf{Term} \\[6pt]
\operatorname{scope}(\mathsf{Num}(n)) &= \mathsf{Par}(\mathsf{Num}(n)) \\[4pt]
\operatorname{scope}(\mathsf{Var}(x)) &= \mathsf{Par}(\mathsf{Var}(x)) \\[4pt]
\operatorname{scope}(\mathsf{Par}(t)) &= \mathsf{Par}(\operatorname{scope}(t)) \\[4pt]
\operatorname{scope}(\mathsf{App}(t,u)) &= \mathsf{Par}(\operatorname{scope}(t)\ \operatorname{scope}(u)) \\[4pt]
\operatorname{scope}(\mathsf{Bin}(o,t,u)) &= \mathsf{Par}(\operatorname{scope}(t)\ o\ \operatorname{scope}(u)) \\[12pt]

t \ \text{is in normal-form scopes}
&\Longleftrightarrow
t = \operatorname{scope}(t)
\end{align}

\section{Appendix S: Spherepop Calculus --- Small-Step Semantics (Inference Rules)}

\begin{align}
t &\to t' \qquad \text{(one-step reduction)} \\[12pt]

\frac{t\to t'}{\mathsf{Par}(t)\to \mathsf{Par}(t')}
\qquad
\frac{t\to t'}{\mathsf{Bin}(o,t,u)\to \mathsf{Bin}(o,t',u)}
\qquad
\frac{u\to u'}{\mathsf{Bin}(o,t,u)\to \mathsf{Bin}(o,t,u')}
\\[16pt]

\frac{t\to t'}{\mathsf{App}(t,u)\to \mathsf{App}(t',u)}
\qquad
\frac{u\to u'}{\mathsf{App}(t,u)\to \mathsf{App}(t,u')}
\\[16pt]

\frac{}{\mathsf{Bin}(+, \mathsf{Num}(a), \mathsf{Num}(b)) \to \mathsf{Num}(a+b)}
\qquad
\frac{}{\mathsf{Bin}(-, \mathsf{Num}(a), \mathsf{Num}(b)) \to \mathsf{Num}(a-b)}
\\[14pt]

\frac{}{\mathsf{Bin}(\times, \mathsf{Num}(a), \mathsf{Num}(b)) \to \mathsf{Num}(ab)}
\qquad
\frac{b\neq 0}{\mathsf{Bin}(\div, \mathsf{Num}(a), \mathsf{Num}(b)) \to \mathsf{Num}(a/b)}
\\[14pt]

\frac{}{\mathsf{Bin}(\wedge, \mathsf{Num}(a), \mathsf{Num}(b)) \to \mathsf{Num}(a^b)}
\\[18pt]

t \Downarrow v
\qquad
\frac{t\to t'\quad t'\Downarrow v}{t\Downarrow v}
\qquad
\frac{}{v\Downarrow v}
\end{align}

\section{Appendix T: Spherepop Calculus --- Normalization and Confluence (Derivation Schemas)}

\begin{align}
\mathcal{R} &= \{(l_i,r_i)\}_{i\in I} \qquad \text{(rewrite system)} \\[10pt]

t \Rightarrow_{\mathcal{R}} t'
&\Longleftrightarrow
\exists\, C[\ ]:\ \exists\,(l,r)\in\mathcal{R}:\ t=C[l],\ t'=C[r]
\\[14pt]

\Rightarrow_{\mathcal{R}}^{\ast}
&=
\text{reflexive-transitive closure}
\\[14pt]

\mathrm{NF}(t)
&=
\{u \mid t\Rightarrow_{\mathcal{R}}^{\ast}u\ \wedge\ \neg\exists u': u\Rightarrow_{\mathcal{R}} u'\}
\\[14pt]

\text{(Local Confluence)}
\qquad
t\Rightarrow_{\mathcal{R}} t_1\ \wedge\ t\Rightarrow_{\mathcal{R}} t_2
&\Longrightarrow
\exists u:\ t_1\Rightarrow_{\mathcal{R}}^{\ast}u\ \wedge\ t_2\Rightarrow_{\mathcal{R}}^{\ast}u
\\[14pt]

\text{(Termination)}
\qquad
\exists\, \prec\ \text{well-founded}:\ 
t\Rightarrow_{\mathcal{R}} t'
&\Longrightarrow
t' \prec t
\\[14pt]

\text{(Newman)}
\qquad
\text{Termination} \ \wedge\ \text{Local Confluence}
&\Longrightarrow
\text{Confluence}
\\[14pt]

\text{(Unique Normal Form)}
\qquad
\text{Confluence} \ \wedge\ \text{Termination}
&\Longrightarrow
\forall t,\ \exists!\ u\in \mathrm{NF}(t)
\\[14pt]

\operatorname{scope}(t) \Rightarrow_{\mathcal{R}}^{\ast} u
&\Longrightarrow
u = \operatorname{scope}(u)
\end{align}

\section{Appendix U: Spherepop Calculus --- Static Semantics (Typing Judgments)}

\begin{align}
\Gamma &:\ \mathsf{VarName} \rightharpoonup \mathsf{Type} \\[6pt]
\tau &\in \mathsf{Type} ::= \mathsf{Nat}\mid \mathsf{Rat}\mid \mathsf{Real} \\[10pt]

\Gamma \vdash t : \tau
\\[14pt]

\frac{}{ \Gamma \vdash \mathsf{Num}(n) : \mathsf{Nat}}
\qquad
\frac{x\in \mathrm{dom}(\Gamma)}{\Gamma \vdash \mathsf{Var}(x) : \Gamma(x)}
\\[16pt]

\frac{\Gamma \vdash t:\tau \quad \Gamma \vdash u:\tau}{\Gamma \vdash \mathsf{App}(t,u):\tau}
\\[16pt]

\frac{\Gamma \vdash t:\mathsf{Nat}\quad \Gamma \vdash u:\mathsf{Nat}}{\Gamma \vdash \mathsf{Bin}(+,t,u):\mathsf{Nat}}
\qquad
\frac{\Gamma \vdash t:\mathsf{Nat}\quad \Gamma \vdash u:\mathsf{Nat}}{\Gamma \vdash \mathsf{Bin}(-,t,u):\mathsf{Nat}}
\\[14pt]

\frac{\Gamma \vdash t:\mathsf{Nat}\quad \Gamma \vdash u:\mathsf{Nat}}{\Gamma \vdash \mathsf{Bin}(\times,t,u):\mathsf{Nat}}
\qquad
\frac{\Gamma \vdash t:\mathsf{Nat}\quad \Gamma \vdash u:\mathsf{Nat}\quad u\neq 0}{\Gamma \vdash \mathsf{Bin}(\div,t,u):\mathsf{Rat}}
\\[14pt]

\frac{\Gamma \vdash t:\mathsf{Nat}\quad \Gamma \vdash u:\mathsf{Nat}}{\Gamma \vdash \mathsf{Bin}(\wedge,t,u):\mathsf{Nat}}
\\[16pt]

\frac{\Gamma \vdash t:\tau}{\Gamma \vdash \mathsf{Par}(t):\tau}
\end{align}

\section{Appendix V: Spherepop Calculus --- Progress and Preservation (Derivation Forms)}

\begin{align}
\mathrm{Val}(t)
&\Longleftrightarrow
t=\mathsf{Num}(n)
\\[12pt]

\text{Preservation:}\qquad
\Gamma \vdash t:\tau \ \wedge\ t\to t'
&\Longrightarrow
\Gamma \vdash t':\tau
\\[14pt]

\text{Progress:}\qquad
\emptyset \vdash t:\tau
&\Longrightarrow
\mathrm{Val}(t)\ \vee\ \exists t': t\to t'
\\[16pt]

\frac{\Gamma \vdash t:\tau \quad t\to t'}{\Gamma \vdash t':\tau}
\qquad
\frac{\emptyset \vdash t:\tau \quad \neg\mathrm{Val}(t)}{\exists t': t\to t'}
\end{align}

\section{Appendix W: Parenthesis-Completion Normal Form --- Totality and Idempotence}

\begin{align}
\operatorname{scope} &: \mathsf{Term}\to \mathsf{Term} \\[10pt]

\forall t\in \mathsf{Term},\quad
\operatorname{scope}(t)
&=
\mathsf{Par}(t')\ \text{for some }t'\in\mathsf{Term}
\\[14pt]

\operatorname{scope}(\operatorname{scope}(t))
&=
\operatorname{scope}(t)
\\[14pt]

\forall t,\quad t=\operatorname{scope}(t)
&\Longleftrightarrow
t\ \text{has explicit scopes at every node}
\\[16pt]

t \equiv_{\alpha} u
&\Longleftrightarrow
\operatorname{erasePar}(t)=\operatorname{erasePar}(u)
\\[12pt]

\operatorname{erasePar}(\mathsf{Par}(t)) &= \operatorname{erasePar}(t) \\[4pt]
\operatorname{erasePar}(\mathsf{Num}(n)) &= \mathsf{Num}(n) \\[4pt]
\operatorname{erasePar}(\mathsf{Var}(x)) &= \mathsf{Var}(x) \\[4pt]
\operatorname{erasePar}(\mathsf{App}(t,u)) &= \mathsf{App}(\operatorname{erasePar}(t),\operatorname{erasePar}(u)) \\[4pt]
\operatorname{erasePar}(\mathsf{Bin}(o,t,u)) &= \mathsf{Bin}(o,\operatorname{erasePar}(t),\operatorname{erasePar}(u))
\\[14pt]

\operatorname{erasePar}(\operatorname{scope}(t))
&=
\operatorname{erasePar}(t)
\\[14pt]

t \Rightarrow_{\mathcal{R}}^{\ast} u
&\Longrightarrow
\operatorname{scope}(t)\Rightarrow_{\mathcal{R}}^{\ast}\operatorname{scope}(u)
\end{align}

\section{Appendix X: Parsing Uniqueness for the Concrete Grammar (Inductive Schemas)}

\begin{align}
\mathsf{parse} &: \Sigma^\ast \rightharpoonup \mathsf{Term} \\[10pt]
w &\in \Sigma^\ast \\[10pt]

\text{Unambiguity:}\qquad
\mathsf{parse}(w)=t_1\ \wedge\ \mathsf{parse}(w)=t_2
&\Longrightarrow
t_1=t_2
\\[14pt]

\text{Structural Induction Schema:}\qquad
\Big(\forall n,\ P(\mathsf{Num}(n))\Big)
\wedge
\Big(\forall x,\ P(\mathsf{Var}(x))\Big)
\wedge
\Big(\forall t,\ P(t)\Rightarrow P(\mathsf{Par}(t))\Big)
\\
\wedge
\Big(\forall t,u,\ P(t)\wedge P(u)\Rightarrow P(\mathsf{App}(t,u))\Big)
\wedge
\Big(\forall o,t,u,\ P(t)\wedge P(u)\Rightarrow P(\mathsf{Bin}(o,t,u))\Big)
&\Longrightarrow
\forall t,\ P(t)
\\[18pt]

\operatorname{prec}(\wedge) &> \operatorname{prec}(\times,\div) > \operatorname{prec}(+,-) \\[10pt]
\operatorname{assoc}(\wedge) &= \text{right} \\[6pt]
\operatorname{assoc}(\times,\div) &= \text{left} \\[6pt]
\operatorname{assoc}(+,-) &= \text{left}
\end{align}

\section{Appendix Y: Spherepop Expressions as Venn and Euler Diagrams}

\begin{align}
\mathcal{U} &: \text{universal set} \\[6pt]
A,B,C &\subseteq \mathcal{U} \\[10pt]

\llbracket t \rrbracket
&:
\mathsf{Term} \to \mathcal{P}(\mathcal{U})
\\[12pt]

\llbracket \mathsf{Num}(n) \rrbracket
&=
U_n \subseteq \mathcal{U}
\\[10pt]

\llbracket \mathsf{Var}(x) \rrbracket
&=
U_x \subseteq \mathcal{U}
\\[12pt]

\llbracket \mathsf{Par}(t) \rrbracket
&=
\llbracket t \rrbracket
\\[12pt]

\llbracket \mathsf{Bin}(+,t,u) \rrbracket
&=
\llbracket t \rrbracket \cup \llbracket u \rrbracket
\\[10pt]

\llbracket \mathsf{Bin}(\times,t,u) \rrbracket
&=
\llbracket t \rrbracket \cap \llbracket u \rrbracket
\\[10pt]

\llbracket \mathsf{Bin}(-,t,u) \rrbracket
&=
\llbracket t \rrbracket \setminus \llbracket u \rrbracket
\\[10pt]

\llbracket \mathsf{Bin}(\div,t,u) \rrbracket
&=
\llbracket t \rrbracket \cap \llbracket u \rrbracket^{c}
\\[12pt]

\llbracket \mathsf{App}(t,u) \rrbracket
&=
\llbracket t \rrbracket \cap \llbracket u \rrbracket
\\[14pt]

\llbracket \operatorname{scope}(t) \rrbracket
&=
\llbracket t \rrbracket
\\[12pt]

\llbracket t \rrbracket
=
\bigcup_{i}
R_i
\qquad
R_i = \bigcap_{j}
U_{ij}
\\[12pt]

\forall t,\quad
\llbracket t \rrbracket
&=
\text{finite Euler decomposition}
\\[14pt]

t_1 \equiv t_2
&\Longleftrightarrow
\llbracket t_1 \rrbracket = \llbracket t_2 \rrbracket
\end{align}

\section{Appendix Z: Spherepop Expressions as Programmable Circuits}

\begin{align}
\mathcal{C} &: \text{symmetric monoidal category} \\[6pt]
\otimes &: \mathcal{C}\times\mathcal{C}\to\mathcal{C} \\[6pt]
I &: \text{monoidal unit} \\[10pt]

\mathsf{Wire}_\tau
&\in \mathrm{Ob}(\mathcal{C})
\\[10pt]

\mathsf{Gate}_o
&:
\mathsf{Wire}_{\tau_1}\otimes \mathsf{Wire}_{\tau_2}
\to
\mathsf{Wire}_{\tau_3}
\\[12pt]

\llbracket\!\llbracket t \rrbracket\!\rrbracket
&:
\mathsf{Wire}_{\tau_1}\otimes\cdots\otimes\mathsf{Wire}_{\tau_n}
\to
\mathsf{Wire}_{\tau}
\\[14pt]

\llbracket\!\llbracket \mathsf{Num}(n) \rrbracket\!\rrbracket
&=
\mathsf{Const}_n : I \to \mathsf{Wire}_{\tau}
\\[12pt]

\llbracket\!\llbracket \mathsf{Var}(x) \rrbracket\!\rrbracket
&=
\mathsf{Id}_{\mathsf{Wire}_{\tau}}
\\[12pt]

\llbracket\!\llbracket \mathsf{Par}(t) \rrbracket\!\rrbracket
&=
\llbracket\!\llbracket t \rrbracket\!\rrbracket
\\[12pt]

\llbracket\!\llbracket \mathsf{Bin}(o,t,u) \rrbracket\!\rrbracket
&=
\mathsf{Gate}_o
\circ
\left(
\llbracket\!\llbracket t \rrbracket\!\rrbracket
\otimes
\llbracket\!\llbracket u \rrbracket\!\rrbracket
\right)
\\[14pt]

\llbracket\!\llbracket \mathsf{App}(t,u) \rrbracket\!\rrbracket
&=
\mathsf{Gate}_{\mathsf{apply}}
\circ
\left(
\llbracket\!\llbracket t \rrbracket\!\rrbracket
\otimes
\llbracket\!\llbracket u \rrbracket\!\rrbracket
\right)
\\[14pt]

\llbracket\!\llbracket t_1;t_2 \rrbracket\!\rrbracket
&=
\llbracket\!\llbracket t_2 \rrbracket\!\rrbracket
\circ
\llbracket\!\llbracket t_1 \rrbracket\!\rrbracket
\\[14pt]

\llbracket\!\llbracket t \rrbracket\!\rrbracket
&=
\text{string diagram in }\mathcal{C}
\\[12pt]

t_1 \equiv t_2
&\Longleftrightarrow
\llbracket\!\llbracket t_1 \rrbracket\!\rrbracket
=
\llbracket\!\llbracket t_2 \rrbracket\!\rrbracket
\end{align}

\newpage
\begin{thebibliography}{99}

\bibitem{Goodhart1975}
C. A. E. Goodhart.
\textit{Problems of Monetary Management: The U.K. Experience}.
In \textit{Papers in Monetary Economics}, Reserve Bank of Australia, 1975.

\bibitem{CampbellKelly2003}
M. Campbell-Kelly.
\textit{From Airline Reservations to Sonic the Hedgehog: A History of the Software Industry}.
MIT Press, 2003.

\bibitem{Norman2013}
D. A. Norman.
\textit{The Design of Everyday Things}.
Revised and Expanded Edition, Basic Books, 2013.

\bibitem{Suchman2007}
L. A. Suchman.
\textit{Human--Machine Reconfigurations: Plans and Situated Actions}.
Second Edition, Cambridge University Press, 2007.

\bibitem{Simon1996}
H. A. Simon.
\textit{The Sciences of the Artificial}.
Third Edition, MIT Press, 1996.

\bibitem{Shannon1948}
C. E. Shannon.
A Mathematical Theory of Communication.
\textit{Bell System Technical Journal}, 27(3):379--423, 1948.

\bibitem{Floridi2011}
L. Floridi.
\textit{The Philosophy of Information}.
Oxford University Press, 2011.

\bibitem{Lessig2006}
L. Lessig.
\textit{Code: Version 2.0}.
Basic Books, 2006.

\bibitem{Keller2013}
E. F. Keller.
\textit{Making Sense of Life: Explaining Biological Development with Models, Metaphors, and Machines}.
Harvard University Press, 2003.

\bibitem{WinogradFlores1986}
T. Winograd and F. Flores.
\textit{Understanding Computers and Cognition: A New Foundation for Design}.
Ablex Publishing, 1986.

\bibitem{OReilly2005}
T. O'Reilly.
What Is Web 2.0: Design Patterns and Business Models for the Next Generation of Software.
\textit{O'Reilly Media}, 2005.

\bibitem{Kittler1999}
F. A. Kittler.
\textit{Gramophone, Film, Typewriter}.
Stanford University Press, 1999.

\bibitem{Raymond1999}
E. S. Raymond.
\textit{The Cathedral and the Bazaar}.
O'Reilly Media, 1999.

\bibitem{ViHart2013}
V. Hart.
Confessions of a Mathematician.
Online video lecture series, 2013.

\bibitem{VimTutor}
B. Moolenaar et al.
\textit{vimtutor}: Interactive Tutorial for the Vim Text Editor.
Included with Vim distributions, various editions.

\bibitem{VimAdventures}
Z. Sklar.
\textit{Vim Adventures}.
Educational video game, online release, 2010.

\bibitem{MicrosoftWindows31}
Microsoft Corporation.
\textit{Microsoft Windows Version 3.1 User’s Guide}.
Microsoft Press, 1992.

\end{thebibliography}

\end{document}
